\section{Рекурсивная модель и переиспользование малых языковых моделей}

\subsection{Общая схема}

Предлагаемая схема состоит из трех частей. Первая часть --- малая предобученная языковая модель, которая преобразует префикс в начальное скрытое состояние и начальное распределение. Вторая часть --- общий оператор рекурсивного уточнения, применяемый несколько раз. Третья часть --- модель времени остановки, которая на каждом шаге оценивает распределение \(\tau^\star\).

На рис.~\ref{fig:recursive_scheme} показана структура вычислений. Диаграмма намеренно оставлена простой: она отражает только те блоки, которые нужны для постановки задачи.


\begin{figure}[H]
\centering
\fbox{%
\begin{minipage}{0.92\textwidth}
\centering
\small
\begin{tabular}{c c c}
Префикс $X_{1:t}$ & $\longrightarrow$ & малая языковая модель \\
 & & $\downarrow$ \\
 & & начальное состояние $Z_{t,0}$ \\
 & & $\downarrow$ \\
рекурсивное уточнение & $\longrightarrow$ & распределение $Q_{t,k}$ \\
$\downarrow$ & & $\circlearrowleft$ продолжить \\
модель времени $\widehat p_{t,k}$ & $\longrightarrow$ & правило остановки
\end{tabular}
\end{minipage}}
\caption{Рекурсивная схема прогнозирования речи и оценки времени остановки}
\label{fig:recursive_scheme}
\end{figure}

Смысл рекурсии состоит не в механическом увеличении числа слоев. Один и тот же оператор уточнения применяется к состоянию несколько раз, а число применений выбирается по данным. Поэтому модель может тратить меньше вычислений на простые позиции и больше --- на сложные.

\subsection{Сравнение с низкоранговой донастройкой}

Низкоранговая донастройка LoRA и предлагаемая рекурсивная схема решают разные задачи. LoRA уменьшает число обучаемых параметров: основные матрицы предобученной модели замораживаются, а изменение весов задается через произведение двух матриц малого ранга. Это статическая донастройка параметров. После обучения вычислительная схема остается фиксированной.

Рекурсивная схема управляет не только параметрами, но и числом внутренних шагов. Она может быть построена поверх той же малой предобученной модели и даже использовать LoRA внутри оператора уточнения. Поэтому сравнение должно быть проведено не как противопоставление, а как проверка трех вариантов: только низкоранговая донастройка, только рекурсивное уточнение, рекурсивное уточнение с низкоранговыми добавками. Таблица~\ref{tab:lora_compare_design} задает минимальный план такого сравнения.

\begin{table}[H]
\centering
\caption{Сравнение способов переиспользования малой языковой модели}
\label{tab:lora_compare_design}
\begin{tabularx}{\textwidth}{lYYY}
\toprule
Вариант & Что обучается & Что измеряется & Ожидаемый смысл сравнения \\
\midrule
Замороженная основа & только модель времени & качество, стоимость & нижняя граница затрат на дообучение \\
Низкоранговая донастройка & добавочные низкоранговые матрицы & качество при фиксированной глубине & статическая параметрическая адаптация \\
Рекурсивное уточнение & общий оператор уточнения и модель времени & качество--стоимость & динамический выбор числа шагов \\
Рекурсия с низкоранговыми добавками & оператор уточнения с малым числом параметров & качество--стоимость, число параметров & совместное использование обеих идей \\
\bottomrule
\end{tabularx}
\end{table}

\subsection{Варианты переиспользования}

Чтобы не обучать языковую модель с нуля, предлагается следующий порядок. Сначала берется малая предобученная модель и используется как исходный преобразователь \(E_\phi\). Ее нижние слои замораживаются, поскольку они кодируют частотные и локальные свойства текста. Затем обучается небольшая проекция скрытого состояния в пространство рекурсивного уточнения. После этого обучаются оператор \(U_\theta\) и модель времени остановки.

Более экономный вариант состоит в заморозке всей исходной модели. Тогда рекурсия работает как надстройка над ее скрытыми состояниями. Такой вариант проще, но может оказаться недостаточно гибким. Более гибкий вариант допускает дообучение верхних слоев или низкоранговые добавки в них. В экспериментальном плане эти варианты должны сравниваться отдельно, иначе нельзя понять, что именно дало выигрыш: предобучение, рекурсия или дополнительные параметры.

